\chapter{Of Marriage, the Lottery, and the Right of Birth}
\label{art:marriage}

\articlenote{In which is set down the marriage bond, the lottery by which the silo permits the bearing of children, and the right of every citizen to petition to be a parent, subject to the silo's need to preserve its population in balance.}

\section{The Marriage Bond}

A marriage is a partnership of two citizens, contracted before Judicial and entered into the rolls under Article 2, by which the two citizens become a household unit and are bound to each other and to the silo in such further duties as this Pact sets down. The marriage bond is the ordinary form by which two citizens declare their commitment to share labor, resources, and, where the lottery permits, the bearing and raising of children. It asks of each party a constancy the silo itself depends upon no less than it depends on breath and water: a household divided against itself cannot be trusted to raise a child in the silo's peace, and a silo built of divided households cannot be trusted to keep its own.

\begin{clauses}
\item A marriage is dissolved only by the death of one of the parties, or by mutual petition granted by Judicial under such terms as Judicial finds just. A marriage dissolved remains entered in the Judicial archive, and does not erase the citizenship rolls of either party.
\item Two citizens may petition to be married before Judicial; Judicial shall not refuse such a petition save where one of the parties is already married or is a ward unable to consent under Article 2.
\item A marriage entered creates no obligation of children; the marriage bond stands as a household association alone, until and unless the partners petition the lottery under Section 2.
\item Housing shall accommodate married couples as married persons under Article 14, allocating housing by the partnership's size and needs rather than counting the two as separate residents.
\end{clauses}

\section{The Lottery for Children}

The silo's population must remain within the bounds set by the silo's stores and the capacity of its Departments to provide for all its citizens. To this end, the bearing of children is governed by a lottery held once yearly: citizens may petition to bear a child, and the lottery selects from among the petitioners such number as the silo's population can sustain.

\begin{clauses}
\item Any citizen, whether married or unmarried, or any partnership of two citizens married under Section 1, may petition to the lottery to be granted the right to bear a child. A petition shall name the petitioner(s) and the intended bearing parent, if different.
\item The lottery is drawn by Judicial, with the Mayor and the Head of Supply present, once yearly. The number of slots granted is set by the Mayor, in consultation with the Head of Supply and the Heads of other Departments, based on the silo's population capacity for the coming year as recorded in the census under Article 2.
\item A citizen granted a slot in the lottery receives the right to bear a child, and the mandatory birth-control measure, previously fitted, is deactivated for such time as permits one conception and bearing. The window for conception shall be no fewer than three seasons and no longer than one full year, at the direction of the Infirmary under Article 19.
\item A citizen who wins the lottery may gift or transfer the right to bear a child to a spouse, a sibling, or another citizen of their household, provided Judicial approves the transfer; such a transfer does not require the recipient to be married or partnered.
\item No citizen shall be compelled to bear a child, and the winning of a lottery slot does not obligate the winner to use it. A slot not used within the window of conception shall be forfeited, and the citizen's birth-control measure shall be re-activated.
\end{clauses}

\section{The Right of Birth and the Window of Conception}

A citizen granted the lottery's right to bear a child must do so within the window of conception granted; outside this window, the child-bearing is unlicensed and unsanctioned. A child born within the window and properly attested shall be named and entered in the rolls under Article 2 as a full citizen.

\begin{clauses}
\item A child conceived within the lottery window and born to a citizen holding the lottery right shall be named and entered in the census by Judicial within ten days of birth, exactly as any other citizen is named.
\item A child born to a citizen who did not win the lottery, or born outside the window of conception granted, is still a citizen and shall be named and entered in the rolls; but the citizen(s) responsible for the child's birth are answerable under Article 17, as set down in Section 3.
\item The Infirmary shall keep, for every citizen of bearing age, a record of the citizen's birth-control status and the dates of any lottery window granted. These records are kept in the Judicial archive under Article 6 and are available to Judicial for verification.
\end{clauses}

\section{Unsanctioned Children and the Answering for Unlicensed Birth}

A child conceived and borne without the lottery's sanction is no less a citizen than any other; the child receives full standing, a name, and entry in the census. But the adult(s) responsible for the child's conception are answerable under Article 17 for the breach of the silo's population law.

\begin{clauses}
\item The citizen(s) responsible for an unsanctioned child's conception shall be determined by Judicial, taking evidence from the child's birth records, the parent(s), and any other source; once responsibility is assigned, the responsible citizen(s) is answerable under Article 17.
\item The bearing of an unsanctioned child does not prevent a citizen from petitioning the lottery in future years; a citizen who has borne an unsanctioned child and is later granted a lottery slot may use it for a sanctioned bearing.
\item No child born outside the lottery shall be treated differently from any other citizen in the silo's law, ration, housing, or any other respect. The penalty applies to the responsible adult(s), not to the child.
\end{clauses}

\section{The Administration of the Lottery}

The lottery is held once yearly, its slots are announced, and the process is kept in the Judicial archive under Article 6. The lottery is a matter of great consequence to every citizen, and transparency is the guard against favoritism or secret manipulation.

\begin{clauses}
\item The lottery shall be drawn in the Assembly under Article 20, or at such place as the Assembly determines, and any citizen may witness the drawing. The names of the winners shall be announced publicly and entered in the Judicial archive.
\item A citizen who believes the lottery was conducted unfairly may petition Judicial for a ruling; Judicial may order the lottery re-drawn if sufficient cause is found.
\item The slots granted by the lottery are the sole means by which a citizen's birth-control measure may be deactivated. No other authority—not the Mayor, not a Department head, not Judicial itself—may order such deactivation.
\item A citizen denied a lottery slot may petition again in following years without penalty; the lottery is held once yearly without exception, and no citizen is permanently barred from petitioning.
\end{clauses}

\section{The Duty to Maintain the Population}

The silo's survival depends on its population remaining stable across the generations. The lottery is the mechanism by which this balance is struck: the silo grows when deaths create new slots, and the silo sustains when births replace the dead. This balance is not a cruelty, but the silo's means of keeping every citizen fed and the works of the silo in motion for generations to come.

It is recorded, though not by the name of the household, that a citizen who waited three years for a lottery slot, and buried a stillborn child in the fourth, came before the Assembly and said that this balance, hard as it had been to bear, was the same balance that had kept her own mother fed as a child, and asked that no citizen be spared the waiting who had not first been spared the alternative. This Pact does not require any citizen to feel as she felt. It asks only that every citizen know such a citizen has lived, and judged as she judged.

\begin{clauses}
\item No citizen is deprived of the chance to bear a child by this Article alone; only the timing is constrained. Citizens who are denied slots in one year may petition in the next, and the lottery is held without fail, every year, for so long as the silo stands.
\item The capacity of the silo to support children is set by the Mayor and the Heads of Departments based on the census; as the silo's stores grow or shrink, as Departments expand or contract, the number of lottery slots may grow or shrink accordingly. The Mayor shall announce the year's slot count before the lottery is drawn.
\end{clauses}
